In a five-seed CPU study on 1D Burgers, a reimplemented FNO is far more accurate than a DeepONet, an MLP and a CNN at 128 points with 400 training pairs and the same 100-epoch budget, and its advantage remains at 100 and 200 training pairs. The DeepONet and CNN baselines are under-fit at that budget, and the DeepONet result is weaker than published comparisons suggest, so the size of the gap should not be generalised. Applied zero-shot to 256 and 512 points, its error is unchanged, whereas the error of the natively applied CNN grows from \rhval{main/cnn/burgers_nu0.01/rel_l2/mean:2} to \rhval{main/cnn/burgers_nu0.01/rel_l2_256/mean:2} and \rhval{main/cnn/burgers_nu0.01/rel_l2_512/mean:2}; the flat errors of MLP and DeepONet come from our subsampling rule and should not be read as invariance. These conclusions are limited to one equation, band-limited data and a tuned-learning-rate-only baseline setup, and the resolution test is not a test of super-resolution beyond the training band.
